\begin{enumerate}
    \item Le bilan des forces~:
          \begin{itemize}
              \item le poids du skieur $\vec{P}$ de direction verticale, de sens
                    vers le bas et d'intensité
                    $P=m\times g=90\times10=\qty{900}{\N}$.
                    \hfill\textbf{(0,25pt)}
              \item la réaction de la piste $\vec{R}$ est décomposée en deux
                    forces~:
                    \begin{itemize}
                        \item $\vec{R}_{N}$~: la réaction normale de la piste,
                              de sens vers le haut,
                              \hfill\textbf{(0,25pt)}
                        \item $\vec{f}$~: la force de frottement exercée par la
                              piste sur le skieur, de sens opposé au mouvement
                              et d'intensité $f=\qty{30}{\N}$.
                              \hfill\textbf{(0,25pt)}
                    \end{itemize}
              \item la tension du câble $\vec{T}$ de direction celle du câble,
                    de sens vers le haut du câble.
                    \hfill\textbf{(0,25pt)}
          \end{itemize}
    \item \mbox{}\par
          \begin{center}
              \begin{tikzpicture}
              \node[anchor=south west,inner sep=0] (image) at (0,0)
                  {\includegraphics[width=0.6\textwidth]{2025_12_01_2444c970f150c248fc83g-3}};
              \begin{scope}[x={(image.south east)},y={(image.north west)}]
                  \draw[blue,ultra thick,->] (0.42,0.49) -- ++(0,-0.5)
                      node[anchor=south west,inner ysep=0pt] {$\vec{P}$};
                  \draw[red,ultra thick,->] (0.446,0.39) -- ++(100:0.5)
                      node[anchor=north east,inner ysep=0pt] {$\vec{R}_{N}$};
                  \draw[red,ultra thick,->] (0.446,0.39) -- ++(219:0.32)
                      node[anchor=south west,inner xsep=0pt] {$\vec{f}$};
                  \draw[Green,ultra thick,->] (0.436,0.49) -- ++(75:0.5)
                      node[anchor=north west,inner ysep=0pt] {$\vec{T}$};
              \end{scope}
              
          \end{tikzpicture}
          \end{center}
    \item $W_{AB}(\vec{P})=\vec{P}\cdot \vv{AB}=
          P\cdot AB\cdot\cos(\alpha+\ang{90})$
          \hfill\textbf{(0,5pt)}

          A.N.
          $W_{AB}(\vec{P})=900\times 1250\times\cos(\ang{110})=
          \qty{-3.85e5}{\J}$
          \hfill\textbf{(0,5pt)}
    \item $W_{AB}(\vec{f})=\vec{f}\cdot \vv{AB}=
          f\cdot AB\cdot\cos(\ang{180})=-f\cdot AB$
          \hfill\textbf{(0,5pt)}

          A.N.
          $W_{AB}(\vec{P})=-30\times 1250=\qty{-3.75e4}{\J}$
          \hfill\textbf{(0,5pt)}
    \item On a~: $\vec{R}=\vec{R}_{N}+\vec{f}$ donc, 
          $W_{AB}(\vec{R})=W_{AB}(\vec{R}_{N})+W_{AB}(\vec{f})$
          \hfill\textbf{(0,5pt)}
          
          $W_{AB}(\vec{R})=0+W_{AB}(\vec{f})=\qty{-3.75e4}{\J}$
          \hfill\textbf{(0,5pt)} \\
          car $\vec{R}_{N}$ est perpendiculaire au déplacement.
    \item $W_{AB}(\vec{T})=\vec{T}\cdot \vv{AB}=
          T\cdot AB\cdot\cos(\beta)$
          \hfill\textbf{(0,5pt)}

          A.N.
          $W_{AB}(\vec{T})=675\times 1250\times\cos(\ang{60})=
          \qty{4.22e5}{\J}$
          \hfill\textbf{(0,5pt)}
\end{enumerate}

